\section{Теоретические сведения}

\subsection{NAT и маршрутизация пакетов}

NAT (Network Address Translation)~--- механизм преобразования сетевых адресов
в IP-пакетах при передаче через шлюз~\cite{iptables_man}. В режиме \textbf{MASQUERADE}
исходящий адрес заменяется адресом внешнего интерфейса шлюза,
что позволяет множеству внутренних узлов выходить в Интернет через один внешний IP-адрес.

Для маршрутизации IPv4-пакетов ядро Linux содержит параметр \texttt{net.ipv4.ip\_forward}.
По умолчанию он выключен (\texttt{= 0}); для работы шлюза его необходимо
включить (\texttt{= 1}) в \texttt{/etc/sysctl.conf}.

\subsection{iptables и netfilter}

\textbf{netfilter}~--- подсистема ядра Linux, выполняющая фильтрацию и трансформацию
пакетов.
\textbf{iptables}~--- пользовательский интерфейс для управления netfilter~\cite{iptables_man}.
Основные цепочки (chains) таблицы \texttt{nat}:
\begin{itemize}
  \item \texttt{PREROUTING}~--- обработка входящих пакетов до принятия решения о маршруте;
  \item \texttt{POSTROUTING}~--- обработка исходящих пакетов перед отправкой;
  \item цель \texttt{MASQUERADE}~--- подстановка внешнего IP-адреса вместо внутреннего.
\end{itemize}

\subsection{DHCP}

DHCP (Dynamic Host Configuration Protocol, RFC~2131)~\cite{rfc2131}~--- протокол
автоматической настройки сетевых параметров клиента (IP-адрес, маска,
шлюз, DNS). Клиент отправляет броадкаст \texttt{DHCPDISCOVER}, сервер отвечает
\texttt{DHCPOFFER}, далее клиент подтверждает запрос (\texttt{DHCPREQUEST}) и получает
\texttt{DHCPACK} с параметрами.

\subsection{Netplan}

Netplan~--- утилита настройки сетевых интерфейсов в Ubuntu~\cite{netplan_docs}.
Конфигурация описывается в YAML-файле \texttt{/etc/netplan/*.yaml} и применяется
командой \texttt{netplan apply}. Параметры языка YAML: отступы только пробелами,
значения --- после двоеточия с пробелом.
